\documentclass[10pt,a4paper]{amsart}
\usepackage[margin=19mm]{geometry}
\usepackage{amsmath,amssymb,mathtools}
\usepackage[scr=rsfs,cal=euler]{mathalfa}
\usepackage{xfrac}
\usepackage[unicode,hidelinks,bookmarks=false]{hyperref}
\newcommand{\ie}{i.e.\ }
\newcommand{\cf}{cf.\ }
\newcommand{\resp}{respectively\ }
\newcommand{\Subsection}{\subsection}
\providecommand{\nobreakdash}{\nobreak\hbox{-}}
\setlength{\emergencystretch}{2em}
\multlinegap0pt
\usepackage[shortlabels]{enumitem}
\usepackage{amssymb}
\usepackage{mathtools}
\usepackage{xcolor}
\newtheorem{definition}{Definition}
\newtheorem{example}{Example}
\newtheorem{lemma}{Lemma}
\newcommand{\MSO}{\mathsf{MSO}}
\newcommand{\WRP}{\mathsf{WRP}}
\newcommand{\Z}{\mathbb Z}
\title{A Computational Obstruction to Swapping Area and Dinv:\\ An Automata-Theoretic View of the $q,t$-Catalan Symmetry}
\author{Jineon Baek, Byung-Hak Hwang, Joonhyun La, Hongseok Yang}
\date{Source: arXiv:2609.05005v1 (CC BY 4.0)}
\begin{document}
\maketitle

\noindent\emph{Source and licence.} A Computational Obstruction to Swapping Area and Dinv:\\ An Automata-Theoretic View of the $q,t$-Catalan Symmetry. Version arXiv:2609.05005v1 (CC BY 4.0), distributed by arXiv under the Creative Commons Attribution 4.0 International licence, http://creativecommons.org/licenses/by/4.0/, as recorded in the arXiv licence metadata for this exact version and shown on the abstract page https://arxiv.org/abs/2609.05005v1. Full licence notice: \texttt{licenses/arxiv-2609.05005-CC-BY-4.0.txt}.\par\smallskip

\noindent\textit{Editorial scope.} Counted unit: the article's lemma lem:two-parameter-presburger (source0.tex lines 3674--3688). The article's complete original proof of it is delivered with it (source0.tex lines 3690--3927, one \texttt{proof} environment of 238 lines). No mathematics is written, repaired or re-derived: the statement and the proof are the article's own lines, in the article's own notation, and no retained line is substituted or re-flowed.  Retained uncounted as the source-local context the proof needs: lines 1111--1122; lines 1128--1138; lines 1143--1149; lines 1166--1185; lines 2394--2394; lines 2520--2537; lines 2612--2628; lines 2785--2803; lines 2898--2912.  Nothing was omitted from any retained span, so the package carries no omission record.  No retained line contains a citation, so the wrapper carries no bibliography and no bibliography entry.  No retained line contains a figure environment, an image-inclusion command or drawing code, so no image is delivered and the package declares no asset dependency.  Editorial formatting only, in addition: an AMS \texttt{amsart} preamble that restates the article's own declarations and macros used by the retained lines, namely the environments \texttt{definition}, \texttt{example}, \texttt{lemma} on separate counters, together with the standard packages those lines need.  The retained environments print in this document as Definition 1, Example 1, Lemma 1 in order of appearance, whereas the article prints the same items with the numbering of its own full document.  The complete native source of the article as distributed in the arXiv e-print for this exact version is bundled unchanged as \texttt{sources/209405/source0.tex}.
\begin{definition}[Deterministic additive rank source]
\label{def:rank-source}
A \emph{$d$-dimensional deterministic additive rank source} over an alphabet
$\Sigma$ is a deterministic finite-state scanner whose transitions carry
weights in $\Z^d$.  Formally, it is a tuple
\[
   A=(Q,q_0,\delta,\omega),
\]
where $Q$ is a finite state set, $q_0\in Q$ the initial state,
$\delta\colon Q\times\Sigma\to Q$ a deterministic transition function, and
$\omega\colon Q\times\Sigma\to\Z^d$ a weight function.
\end{definition}

\begin{definition}[Prefix rank]\label{def:prefix-rank}
Let $A$ be a deterministic additive rank source and let
$w=a_1\cdots a_n$.  Define
$q_i=\delta(q_{i-1},a_i)$ for $i=1,\ldots,n$.  The \emph{prefix rank} of $A$
before position $i\in\{1,\ldots,n+1\}$ is the total weight accumulated before
reading position $i$:
\[
   \rho_A^w(i)=\sum_{j<i}\omega(q_{j-1},a_j)\in\Z^d.
\]
We also write $q^A_i:=q_{i-1}$ for the state of $A$ just before position $i$.
\end{definition}

\begin{example}[Height as a rank source]\label{ex:height-rank}
Let $A$ be the one-state source with
$\omega(\bullet,U)=+1$ and $\omega(\bullet,D)=-1$.  Then, $\rho_A^w(i)$ is the
height of $w$ just before position $i$.  In particular, if $P$ is a Dyck path
and $p_j$ is the position of its $j$-th up-step, then
$\rho_A^P(p_j)=a_j$, the $j$-th area-sequence entry.
\end{example}

\begin{definition}[Prefix-additive rank function]
\label{def:prefix-additive-rank}
Fix $k\ge 1$.  For each coordinate $r\in\{1,\ldots,k\}$, fix a
$d$-dimensional deterministic additive rank source
$A_r=(Q_r,q_{0,r},\delta_r,\omega_r)$ on $\Sigma$, together with a fixed
local-correction table
$\beta_r\colon Q_r\times\Sigma\to\Z^d$.  A
\emph{$d$-dimensional prefix-additive rank function} on $k$-tuples of input
positions is a function of the form
\[
   \kappa^w(x_1,\ldots,x_k)
   =c_0+\sum_{r=1}^{k}
       \Bigl(\rho_{A_r}^w(x_r)
             +\beta_r(q^{A_r}_{x_r},a_{x_r})\Bigr),
   \qquad c_0\in\Z^d.
\]
We usually suppress the superscript $w$.  A coordinate that makes no
contribution uses the one-state zero-weight source and the zero local
correction.
\end{definition}

\section[Regular-slice semilinearity for WRP]{Regular-slice semilinearity for $\WRP$}\label{sec:slice-semilinearity}

\begin{example}[Coordinates and height for $W_n$]
\label{ex:slice-coordinates}
Take the height rank source of Example~\ref{ex:height-rank}.  On
$W_n=U(UD)^nD$, its prefix ranks (Definition~\ref{def:prefix-rank}) before the
four kinds of positions are
\[
   \rho(\mathsf u,1,0)=0,\qquad
   \rho(\mathsf v,1,j)=1,\qquad
   \rho(\mathsf v,2,j)=2,\qquad
   \rho(\mathsf z,1,0)=1.
\]
The two values inside the repeated part do not depend on $j$ because each
block $UD$ has total height change zero.  Thus, for this particular rank source
and this particular family, the tag and offset determine the prefix rank even
though the repetition index may grow with $n$.  This constancy is special to
the height source on $W_n$; prefix ranks need not generally be constant along
a regular slice.
\end{example}

\begin{lemma}[MSO conditions on a regular slice]
\label{lem:one-loop-finite-state}
Fix a regular slice $(w_n)_{n\ge0}$ with $w_n=uv^nz$ and an $\MSO$ formula
$\varphi(x_1,\ldots,x_k)$.  For each variable $x_\ell$, fix its region
$\tau_\ell\in\{\mathsf u,\mathsf v,\mathsf z\}$: the prefix $u$, a copy of the
repeated block $v$, or the suffix $z$.  Also fix an offset $s_\ell$ valid in
that region, as in the position encoding preceding
Example~\ref{ex:slice-coordinates}.  If $\tau_\ell=\mathsf v$, let
$j_\ell\in\{1,\ldots,n\}$ be the repetition index; otherwise set $j_\ell=0$.
Let $i_\ell$ be the position encoded by $(\tau_\ell,s_\ell,j_\ell)$.  Then the
set
\[
   \{(n,j_1,\ldots,j_k):
      w_n\models\varphi(i_1,\ldots,i_k)\}
\]
is Presburger-definable.
\end{lemma}

\begin{lemma}[Prefix-additive ranks are piecewise affine on a regular slice]
\label{lem:one-loop-rank-affine}
Fix a regular slice $(w_n)_{n\ge0}$ with $w_n=uv^nz$ and a $d$-dimensional prefix-additive rank
function $\kappa$ on $k$-tuples of positions.  For each coordinate
$\ell\in\{1,\ldots,k\}$, fix a region
$\tau_\ell\in\{\mathsf u,\mathsf v,\mathsf z\}$ and an offset $s_\ell$ valid in
that region.  If $\tau_\ell=\mathsf v$, let
$j_\ell\in\{1,\ldots,n\}$ be the repetition index; otherwise set $j_\ell=0$.
Let $i_\ell$ be the position encoded by $(\tau_\ell,s_\ell,j_\ell)$, and let
$J=\{\ell:\tau_\ell=\mathsf v\}$.  Then, the valid parameter tuples
$(n,(j_\ell)_{\ell\in J})$ admit a finite Presburger-definable partition such
that, on each part,
\[
   \kappa^{w_n}(i_1,\ldots,i_k)
   =b_0+b_n n+\sum_{\ell\in J} b_\ell j_\ell,
\]
where $b_0,b_n,b_\ell\in\mathbb Q^d$ are fixed on that part and the expression
takes values in $\Z^d$ on the part.
\end{lemma}

\begin{lemma}[WRP data on a regular slice]
\label{lem:one-loop-presburger}
Fix a $\WRP$ presentation and a regular slice $(w_n)_{n\ge0}$ with
$w_n=uv^nz$.  For every potential
atom under consideration, fix its copy name and the region and offset of each
coordinate.  Each of the following is then Presburger-definable in $n$ and the
remaining repetition indices:
\begin{enumerate}[label=(\alph*)]
\item whether a given potential atom is selected;
\item whether a given selected atom has a specified output label;
\item whether the first of two selected atoms has the smaller atom rank;
\item whether the first of two selected atoms of equal atom rank comes first in
the tie-order.
\end{enumerate}
\end{lemma}

\begin{lemma}[Arithmetic on the two-parameter family]
\label{lem:two-parameter-presburger}
Fix a $\WRP$ presentation and consider the family
$\{W_{m,n}:m\ge 1,\ n\ge 0\}$.  For every potential atom under consideration,
fix its copy name and the region and offset of each coordinate.  Each of the
following is then Presburger-definable in $m,n$ and the remaining repetition
indices:
\begin{enumerate}[label=(\alph*)]
\item whether a given potential atom is selected;
\item whether a given selected atom has a specified output label;
\item whether the first of two selected atoms has the smaller atom rank;
\item whether the first of two selected atoms of equal atom rank comes first in
the tie-order.
\end{enumerate}
\end{lemma}

\begin{proof}
This is the two-parameter analogue of
Lemma~\ref{lem:one-loop-presburger}, and the proof follows the regular-slice
arguments of Lemmas~\ref{lem:one-loop-finite-state}
and~\ref{lem:one-loop-rank-affine}; the only change is that $W_{m,n}$ has
\emph{three} repeated stretches and \emph{two} unbounded parameters $m,n$, in
place of the single block $v$ and parameter $n$ there.  We use the position
encoding fixed just above.

\vspace{0.5em}
\noindent\emph{Overview.}  Each path $W_{m,n}$ is made of long runs of
identical steps, governed by just the two numbers $m$ and $n$.  Whatever a finite
automaton computes while scanning such a run eventually falls into a periodic
pattern, and any running integer total it keeps grows at a constant rate.  So
every decision a $\WRP$ presentation makes on $W_{m,n}$ depends on $m$, $n$, and
the repetition indices of the atoms involved only through linear expressions and
congruences.  In the language introduced in
Section~\ref{sec:slice-semilinearity}, these decisions are
Presburger-definable.  Turning
that intuition into the lemma's four conditions (a)--(d) is the whole proof.
These conditions are of two kinds: the finite-state conditions (a), (b), (d) ask
whether an $\MSO$ formula holds at the chosen positions, a question about
automaton acceptance, while the atom-rank condition (c) asks which of two atom
ranks is lexicographically smaller.  We first isolate the ingredient shared by
both kinds, the eventual
periodicity of an automaton along a repeated block, and then treat the two kinds
in turn, as parts (a),(b),(d) and part (c) below.

\vspace{0.5em}
\noindent\emph{Three block maps, each eventually periodic.}
We use the following repetition fact for both kinds.  The two parts below
scan $W_{m,n}$ with different machines: the $\MSO$ parts
(a), (b), (d) use a DFA over the \emph{marked} alphabet
$\Sigma\times\{0,1\}^k$ (the marking construction recalled in
Section~\ref{sec:slice-semilinearity}), while the atom-rank part~(c) uses an additive
rank source over the plain alphabet $\Sigma$.  In either case, let $Q$ be its
state set, $\delta^*$ its transition extension, and $\delta_y(q)=\delta^*(q,y)$
the state reached on reading $y$ from $q$.  For the marked DFA, call an
occurrence of a repeated block \emph{marked} if at least one chosen position
lies in it, equivalently if some marker bit on that occurrence is $1$; call it
\emph{unmarked} otherwise.  Reading one unmarked occurrence of each repeated
block induces three fixed transformations of $Q$.  Every letter in such an
occurrence carries marker bits $0$, so $\delta_U$ abbreviates reading the
marked letter $(U,0^k)$, and likewise for $\delta_{UD}$ and $\delta_D$.  For
the rank source, each subscript simply denotes the corresponding plain block.
The three transformations are
\[
   g_{\mathsf i}=\delta_U,\qquad g_{\mathsf m}=\delta_{UD},\qquad
   g_{\mathsf f}=\delta_D,
\]
one per region $\tau\in\{\mathsf i,\mathsf m,\mathsf f\}$, so reading $L$
consecutive unmarked block occurrences in a stretch applies the single map
$g_\bullet^{\,L}$.  There are at most $k$ marked block occurrences in the
$\MSO$ run.  Their transformations need not be powers of $g_\bullet$; as in
Lemma~\ref{lem:one-loop-finite-state}, they are absorbed into the fixed
transformations of the factorisation in parts (a), (b), (d) below.  Each
$g_\bullet$ lies in the finite monoid $Q^Q$, hence its powers are
eventually periodic; taking the larger threshold and a common multiple of the
periods gives a single $k_0$ and $p$ with $g_\bullet^{\,L+p}=g_\bullet^{\,L}$
for $L\ge k_0$ and all three maps.  The reduced exponent
$\widehat L\in\{0,\ldots,k_0+p-1\}$ of Lemma~\ref{lem:one-loop-finite-state}
then satisfies $g_\bullet^{\,L}=g_\bullet^{\,\widehat L}$, and ``$\widehat L=c$''
is a Presburger-definable condition on $L$.  This eventual periodicity, captured by the
Presburger-definable reduced exponent, is the only way the machines' state
behaviour along the repeated stretches enters either part below.

\vspace{0.5em}
\noindent\emph{Parts (a), (b), (d): the finite-state conditions.}
Selection, labelling, and the tie-order are each an $\MSO$ formula
$\varphi(x_1,\ldots,x_k)$ with $k$ free position variables, and we must show that
the tuples of parameters for which $\varphi$ holds at the chosen positions form a
Presburger-definable set.  The marked DFA $M$ over $\Sigma\times\{0,1\}^k$ of the setup
above accepts, by the marking construction, exactly when
$W_{m,n}\models\varphi(i_1,\ldots,i_k)$: it reads $W_{m,n}$ with the letters at
the chosen positions $i_1,\ldots,i_k$ flagged by the $k$ marker bits.  Write
$\widetilde W_{m,n}$ for this \emph{marked word}.  On this word, the state
reached by $M$ from its start state $q_0$ is
$\delta_{\widetilde W_{m,n}}(q_0)$.  The formula $\varphi$ holds at the chosen
positions if and only if this state belongs to $F$.  By the encoding, each marked
variable $x_\ell$ has a fixed region $\tau_\ell$ and offset $s_\ell$ and an
unbounded repetition index $j_\ell$ (in $\{1,\ldots,m\}$, $\{1,\ldots,n\}$, or
$\{1,\ldots,m\}$ according to $\tau_\ell$), and its position $i_\ell$ is the
affine form in $j_\ell,m,n$ recorded there; the unbounded data is the tuple
$(m,n,(j_\ell)_\ell)$.

The goal is now concrete: show that the tuples $(m,n,(j_\ell)_\ell)$ with
$\delta_{\widetilde W_{m,n}}(q_0)\in F$ form a Presburger-definable set.  We
proceed in three steps.  First, for every pair of variables assigned to the same
stretch, we distinguish the cases $j_\ell<j_{\ell'}$,
$j_\ell=j_{\ell'}$, and $j_\ell>j_{\ell'}$.  These comparisons determine
whether the two chosen positions lie in the same occurrence of the repeated
block and, if not, which occurrence is read first.  Second, in each resulting
case, we factor the marked word into the block occurrences containing chosen
positions and the intervening runs of unmarked block occurrences.  The
transformations associated with the marked occurrences are fixed, while only
the lengths of the unmarked runs vary.  Third, for each intervening run of $L$
unmarked block occurrences, we subdivide the case according to the bounded
value $\widehat L$ defined above, which satisfies
$g_\bullet^{\,L}=g_\bullet^{\,\widehat L}$.  After this subdivision, the
composite transition, and hence the truth of $\varphi$, is fixed on each
subcase.

\vspace{0.5em}
\noindent\emph{Splitting into cells.}  Recall that the regions $\tau_\ell$ and
offsets $s_\ell$ are finite tags fixed in advance, so the only unbounded data is
$(m,n,(j_\ell)_\ell)$.  For each pair of variables in the same stretch, split
the parameter space according to which of $j_\ell<j_{\ell'}$,
$j_\ell=j_{\ell'}$, and $j_\ell>j_{\ell'}$ holds.  For example, if two
variables lie in the middle stretch $(UD)^n$, equality means that their chosen
positions lie in the same occurrence of $UD$, while $j_\ell<j_{\ell'}$ means
that the occurrence containing $x_\ell$ is read before the occurrence
containing $x_{\ell'}$.  The order between different stretches is already
fixed: the initial stretch $U^m$ is read first, then $(UD)^n$, and then $D^m$.
Each consistent choice of the pairwise comparisons is given by a conjunction
of linear (in)equalities and therefore defines a Presburger-definable cell.  On
such a cell, equal repetition indices identify the variables whose marker bits
occur in the same block occurrence, and strict inequalities fix the order of
the marked block occurrences.  Because the offsets are also fixed, we know
exactly which marker bits are attached to each letter in every such block
occurrence.  The cells cover the whole parameter space, so it suffices to show
that the accepting tuples inside one cell $C$ form a Presburger-definable set.

\vspace{0.5em}
\noindent\emph{Factoring the run on $C$.}  By the previous step, the order of
the block occurrences containing chosen positions and the marker bits attached
to their letters are fixed on $C$.  Reading $\widetilde W_{m,n}$ from left to right therefore alternates
between maximal runs of unmarked block occurrences and individual marked block
occurrences.  A run of $L$ unmarked occurrences of one block applies the
corresponding power $g_{\mathsf i}^{\,L}$, $g_{\mathsf m}^{\,L}$, or
$g_{\mathsf f}^{\,L}$ of a block map.  A single marked block occurrence applies
one fixed transformation $M\in Q^Q$: the cell determines which chosen
positions lie in that occurrence, and their offsets $s_\ell$ are fixed, so the
letters and marker bits in the occurrence are completely determined.  Writing
$t\le k$ for the number of marked block occurrences, composing these maps in
the order the input is scanned
gives
\[
   \delta_{\widetilde W_{m,n}}
   = g_\bullet^{\,L_t}\circ M_t\circ\cdots\circ M_1\circ g_\bullet^{\,L_0}
\]
(later factors on the left, since reading a concatenation composes the maps),
where $M_1,\ldots,M_t$ are the fixed transformations associated with the
marked block occurrences and
$g_\bullet^{\,L_0},\ldots,g_\bullet^{\,L_t}$ are the unmarked runs surrounding
them; any of the exponents $L$ may be $0$.  Each exponent $L$ is the length of one such run, a
nonnegative linear form in $m,n$ and the repetition indices: within each stretch, the
runs are the gaps before, between, and after its marked block occurrences, summing to the
stretch length ($m$ for the initial and final stretches, $n$ for the middle)
minus the number of marked block occurrences in it.

\vspace{0.5em}
\noindent\emph{Making the run constant.}  Refine $C$ by the reduced exponent $\widehat L$
of each of these finitely many run-lengths, adjoining the Presburger-definable condition
``$\widehat L=c$'' for each.  This splits $C$ into finitely many sub-cells on
each of which every $g_\bullet^{\,L}=g_\bullet^{\,\widehat L}$ is one fixed power
of a block map, so the composite above is a single fixed map $\sigma\in Q^Q$.
Then, $\delta_{\widetilde W_{m,n}}(q_0)=\sigma(q_0)$ is constant on the sub-cell,
and the acceptance test $\sigma(q_0)\in F$ has the same answer
everywhere on it.  The accepting tuples are thus the union of the sub-cells on
which $\sigma(q_0)\in F$, a finite union of Presburger-definable sets.  This settles (a),
(b), and (d).

\vspace{0.5em}
\noindent\emph{Part (c): atom-rank comparison.}
The aim is to decide which of two atoms receives the smaller atom rank.  The
point is that on $W_{m,n}$, every atom rank is an affine vector-valued function
of $m$, $n$, and the
repetition indices, once the parameters are split into finitely many cases by threshold
tests and residues; comparing two affine forms is then expressible 
by a Presburger formula.

\vspace{0.5em}
\noindent\emph{One coordinate source at a time.}  Recall
(Definitions~\ref{def:rank-source} and~\ref{def:prefix-rank})
that a rank source is a finite-state scanner that, as it reads each letter $a$
in a state $q$, also adds an integer vector $\omega(q,a)\in\Z^d$ and reports
the running total;
the \emph{prefix rank} before a position is that total accumulated over the
letters strictly before it, the abstract form of the height.  A
prefix-additive rank function
(Definition~\ref{def:prefix-additive-rank}) adds one such contribution, with a
bounded local correction, for each tuple coordinate.  It is therefore enough
to take one coordinate source $A=(Q,q_0,\delta,\omega)$ and show that its
prefix rank before a position is affine in $m$, $n$, and that position's
repetition index, after a finite split into Presburger-definable cases; the finitely many coordinate
contributions can then be summed.

\vspace{0.5em}
\noindent\emph{Weight accumulated along a stretch is affine.}  Reading one occurrence of a block
from a state $q$ moves $q$ by the corresponding block map ($g_{\mathsf i}$,
$g_{\mathsf m}$, or $g_{\mathsf f}$) and adds a weight $\Omega_\tau(q)\in\Z^d$, the
total of $\omega$ over that block occurrence.  So if a stretch is entered in a state $q$, then
after $\nu$ completed block occurrences, the state is
$g_\bullet^{\,\nu}(q)$, which is eventually periodic in $\nu$: there are a threshold and a period so that, once $\nu$ passes the
threshold, the entering state depends only on $\nu$ modulo the period, while the
finitely many $\nu$ below the threshold are separate.  Call each possibility a
\emph{case}; it is either a value below the threshold or a residue modulo the
period above it.  Fix one such case.  Each completed period then adds the same
fixed weight increment, and the bounded leftover is fixed.  Hence the weight
of the first $\nu$ block occurrences is an affine function of $\nu$ on the
case, with rational coefficients and integer values.  (This is
the additive counterpart of the eventual periodicity used for the $\MSO$
parts: the states cycle, and the running weight grows by a fixed amount per
cycle.)

\vspace{0.5em}
\noindent\emph{Prefix rank by stretch.}  Applying this to each stretch, the prefix rank
before a position $(\tau,s,j)$ is, after splitting $m,n,j$ into finitely many such
cases:
\begin{itemize}
\item $\tau=\mathsf i$, the $j$th block occurrence: the prefix is $U^{j-1}$, so the prefix rank is the
weight of $j-1$ occurrences of $U$ read from $q_0$, affine in $j$ on each case of $j$
(a residue modulo the period of $g_{\mathsf i}$, above a threshold);
\item $\tau=\mathsf m$, offset $s$ in the $j$th block occurrence: the prefix is $U^m$, then
$(UD)^{j-1}$, then the fixed length-$(s-1)$ start of $UD$; the $U^m$ part is affine
in $m$, the $(UD)^{j-1}$ part affine in $j$ (read from the state after $U^m$, which
is fixed on the case of $m$), and the within-block part depends only on $s$
and that entering state, so is bounded and fixed on the case; the prefix rank
is thus affine in $m$ and $j$;
\item $\tau=\mathsf f$, the $j$th block occurrence: the prefix is $U^m$, then $(UD)^n$, then
$D^{j-1}$, whose three completed stretches contribute weights affine in $m$, in
$n$, and in $j$, so the prefix rank is affine in $m,n,j$.
\end{itemize}
The bounded local corrections of Definition~\ref{def:prefix-additive-rank} are
likewise constant on each such case, depending only on the fixed offset,
letter, and entering state.  Summing the finitely many coordinate contributions
and the fixed constant $c_0$ as in that definition, the atom rank is one affine
form in $m,n$ and the relevant repetition indices on each case.

\vspace{0.5em}
\noindent\emph{Comparing two atom ranks.}  For two atoms, refine to a common case; both
atom ranks are then affine vectors in $m,n$ and the repetition indices $j_\ell$.
Whether one is lexicographically smaller than the other is a Boolean
combination of affine equalities and inequalities and the congruences defining
the case.  Clearing the fixed denominators makes all coefficients integral, so
the comparison is expressible by a Presburger formula.  This
settles (c).
\end{proof}

\end{document}
